% Additive accounting identities for Sparse Readout Prism.

\section{SRP Score Accounting Identities}
\label{app:score-accounting}
\label{app:interpreting-sparse-readout-prism-figures}
\label{app:additional-diagnostics}

This appendix states the accounting identities used by the figures and control
tables and fixes the bookkeeping conventions used when a trained readout SAE is
displayed in figures. The identities are the same local score account introduced in
\Cref{sec:method-srp}. All reported scores are evaluated in the
model's raw readout space. If training used centering or row normalization, the
SAE coefficients, reconstructed rows, and residual terms are first mapped back
to that space.

\subsection{Raw-Space Row Accounting}
\label{app:score-accounting-row-factorization}

Let \(w_t\in\R^d\) be the unembedding row for token \(t\). Under the row
preprocessing convention used by the reported analyses,
\[
    \begin{aligned}
    x_t&=\frac{w_t-\mu}{a_t}, \qquad a_t>0,\\
    x_t&=\sum_i \widetilde z_{t,i}\widetilde d_i+\widetilde r_t,
    \end{aligned}
\]
with \(a_t=1\) for unnormalized rows. The raw-space accounting then uses
\[
    z_{t,i}=a_t\widetilde z_{t,i},
    \qquad
    d_i=\widetilde d_i,
    \qquad
    r_t=a_t\widetilde r_t,
\]
so that
\[
    w_t=\mu+\sum_i z_{t,i}d_i+r_t .
\]
The displayed local score account contains only this shared offset, SAE feature
terms, and the explicit residual term; no hidden dense base term is added.

\subsection{Scalar Target Identity}
\label{app:score-accounting-linear-query}

A scalar readout target is a linear function of vocabulary rows,
\[
    q_\alpha = \sum_{t\in\mathcal{V}} \alpha_t w_t,
    \qquad
    s_\alpha(h)=h^\top q_\alpha .
\]
Define the target coefficient and SAE feature contribution as
\[
    \beta_i(\alpha)=\sum_t \alpha_t z_{t,i},
    \qquad
    c_i(h,\alpha)=\beta_i(\alpha)\,h^\top d_i .
\]
Substituting the raw-space row accounting gives the exact grouped identity
\[
\begin{aligned}
    s_\alpha(h)
    &=
    s_{\mu}(h,\alpha)
    +s_{\mathrm{feat}}(h,\alpha)\\
    &\quad
    +s_{\mathrm{resid}}(h,\alpha),\\
    s_{\mu}(h,\alpha)
    &=
    \left(\sum_t \alpha_t\right)h^\top\mu,\\
    s_{\mathrm{feat}}(h,\alpha)
    &=
    \sum_i c_i(h,\alpha),\\
    s_{\mathrm{resid}}(h,\alpha)
    &=
    h^\top\sum_t \alpha_t r_t .
\end{aligned}
\]
For zero-sum contrasts, including pairwise logit differences, group contrasts,
reference contrasts, and top-competitor margins, the offset term
cancels.

\subsection{Error, Sign, and Display Rules}
\label{app:score-accounting-fidelity}
\label{app:display-taxonomy}
\label{app:display-residual-accounting}

For a scalar readout target \(q_\alpha\), the exact and reconstructed scores are
\[
\begin{aligned}
    s_{\mathrm{exact}}
    &=
    h^\top q_\alpha,\\
    s_{\mathrm{recon}}
    &=
    s_{\mu}(h,\alpha)
    +
    s_{\mathrm{feat}}(h,\alpha).
\end{aligned}
\]
The residual score is the reconstruction error,
\[
    \epsilon(h,\alpha)
    =
    s_{\mathrm{exact}}-s_{\mathrm{recon}}
    =
    s_{\mathrm{resid}}(h,\alpha),
\]
and aggregate summaries use
\[
    \rho(h,\alpha)
    =
    \frac{|\epsilon(h,\alpha)|}{|s_{\mathrm{exact}}|+\delta},
\]
with \(\delta=0.5\) logit. The floor prevents arbitrarily small margins from
dominating median error summaries; near the decision boundary, sign agreement and
margin-binned flip rates remain the primary reliability checks. For signed
contrasts, side-specific feature support is evidence-bearing when
\[
    \operatorname{sign}(s_{\mathrm{exact}})
    =
    \operatorname{sign}(s_{\mathrm{recon}}).
\]
Each interpreted figure therefore states the selected target, exact score,
reconstructed score, reconstruction error, relative reconstruction error, and
sign status for contrasts. If a plot displays only the largest bars, the
reported feature sum and residual are still computed from the full active
feature sum. Token-level figures also report tokenization or row-level audit
information where it affects interpretation; brittle single-token contrasts are
handled as explicit group contrasts when needed.
